\begin{pastexam}[label={px:ncu-115-dsa-02}]{中央資工所 115 學年度 第 2 題}
（多選；每題 5 分，每個選項單獨計分，答錯每個選項倒扣 1 分）

Which of the following statements about amortized analysis are true?
\begin{enumerate}[label=\Alph*.]
\item Amortized analysis provides an upper bound on the average cost per operation over any sequence of operations, not assuming any probability distribution.
\item If the amortized cost of an operation is $O(1)$, then every individual operation must also run in $O(1)$ time.
\item The potential method assigns a potential value to the data structure state, which is used to account for future expensive operations.
\item In the accounting method, some operations may be charged more than their actual cost to pay for later operations.
\item Amortized analysis and average-case analysis are equivalent concepts.
\end{enumerate}
\end{pastexam}

\begin{pastsoln}
答案：A、C、D。

\textbf{A 對}。 平攤分析保證的是：\textbf{任何}一個長度 $n$ 的操作序列，總成本都不超過 $n$ 乘上平攤
成本（聚集法就是直接求最壞序列的總成本 $T(n)$，再除以 $n$，資結書第~3.1~節）。 這是最壞情況的
保證，沒有假設輸入的機率分布。

\textbf{B 錯}。 平攤 $O(1)$ 只說平均下來便宜，個別操作可以很貴。 反例是資結書第~3.4~節的動態陣列：
插入的平攤成本是 $O(1)$，但陣列滿了要擴充的那一次，必須把 $k$ 個元素全部搬到新陣列，那一次是
$\Theta(k)$。 二進位計數器的 \proc{Increment} 也一樣，平攤 $O(1)$，但從 $0111\cdots1$ 加一要翻
$k$ 個位元。

\textbf{C 對}。 位能法（資結書第~3.3~節）定義位能函數 $\Phi$，把資料結構的狀態對應到一個數，
第 $i$ 個操作的平攤成本是 $\hat c_i = c_i + \Phi(D_i) - \Phi(D_{i-1})$。 便宜的操作讓位能上升，
等於先存起來；昂貴的操作讓位能下降，用存下來的位能付掉多出來的實際成本。 只要
$\Phi(D_n) \ge \Phi(D_0)$，平攤成本的總和就是實際總成本的上界。

\textbf{D 對}。 記帳法（資結書第~3.2~節）就是讓某些操作多收費，多收的部分當作存款放在資料結構
裡，留給以後的操作用。 例如二進位計數器把位元設成 $1$ 時收 $2$：$1$ 付這次的翻轉，$1$ 存在這個位元
上，以後它被翻回 $0$ 時就用這筆存款付，所以 \proc{Increment} 的平攤成本是 $2$。

\textbf{E 錯}。 平均情況分析要先假設輸入的分布（例如每種排列機率相同），算的是期望值；某些輸入
仍然可能很慢。 平攤分析不談機率，對最壞的操作序列也成立。 例如快速排序的平均情況是
$O(n \log n)$，但特定輸入要 $\Theta(n^2)$；動態陣列連續 $n$ 次插入，不論怎麼插，總成本都是
$O(n)$。 兩者是不同的概念。
\end{pastsoln}
